% =============================================================================
% Introduction -- evidence-hierarchy rewrite
% Drop-in include: % =============================================================================
% Introduction -- evidence-hierarchy rewrite
% Drop-in include: % =============================================================================
% Introduction -- evidence-hierarchy rewrite
% Drop-in include: % =============================================================================
% Introduction -- evidence-hierarchy rewrite
% Drop-in include: \input{sections/intro}
% =============================================================================
\section{Introduction}
\label{sec:intro}

This paper separates three kinds of evidence before making any claim about
reinforcement-learning post-training.  Held-out evaluations are the only basis
we use for capability claims.  Online rewards are evidence about optimization
inside a particular reward environment.  Proxy rewards, partial harnesses, and
stack-specific comparisons are treated as hypotheses or diagnostics unless they
are paired with matched held-out evaluation.  This hierarchy is necessary
because PPO~\citep{schulman2017proximal}, GRPO~\citep{shao2024deepseekmath},
and DPO~\citep{rafailov2024direct} are often discussed as algorithm-level
interventions, while LLM post-training outcomes also depend on base model,
instruction tuning, reward parser, group construction, sampler, backend, LoRA
configuration, checkpoint selection, and evaluation harness.  The resulting
identification problem is the LLM post-training analogue of the reproducibility
concerns identified by Henderson et al.~\citep{henderson2018deep} for deep RL.

On the strongest evidence tier, this paper does \emph{not} show that GRPO
broadly improves reasoning.  The clean held-out capability check is a GSM8K
200-example slice that was not used during training, and it improves only from
$82.0\%$ to $83.3\%$ ($p=0.26$).  We therefore do not claim a significant
held-out math effect.  This negative result determines the scope of the rest of
the paper: training reward curves and proxy environments may explain when a
post-training loop has signal, but they do not establish general reasoning or
tool-use capability.

Within that scope, \ours{} is a systems audit of a group-relative RL runner
rather than an algorithm leaderboard.  The corpus contains 79 runs across
structured tool-call emission, GSM8K, MATH-style tasks, HumanEval-style binary
verification, small and frontier-scale model families, and several training
backends.  The runner is inspired by GRPO, but it is not a faithful reproduction
of canonical GRPO: it keeps group-normalized advantages and a critic-free
structure, but it does not implement the PPO probability ratio and clip of
\citet{schulman2017proximal}, does not keep a frozen reference policy for KL
regularization, performs one optimizer step per rollout batch, and does not use
a completion-only token mask.  A P1-A diagnostic on Qwen3.5-4B measures
$61.6\%$--$89.6\%$ of the full-sequence loss magnitude as coming from prompt
tokens rather than completion tokens.  Claims below therefore concern this
runner and its reward environments, not a universal property of canonical GRPO.

The paper's main positive claim sits one tier below held-out capability: the
runner behaves as a reward-contrast amplifier.  In a critic-free group-relative
update, a prompt is useful only when the sampled group contains distinguishable
reward outcomes.  Under binary or verifier-like rewards, if the current policy
succeeds on prompt $x$ with probability $p_x$ and the group size is $G$, the
probability of obtaining a mixed, gradient-carrying group is approximated by
\[
P(\text{usable}) \approx \frac{1}{N}\sum_x \left[1 - (1-p_x)^G - p_x^G\right].
\]
This mechanism explains why warm-starts, group size, reward density, and task
difficulty matter more than task labels.  The familiar $p_x > 1/G$ rule is only
an expected-one-success heuristic; it is not a threshold theorem for learning.

Zero-Variance Fraction (ZVF) and Gradient Utilization (GU) operationalize that
mechanism as diagnostics, not as capability metrics.  ZVF is the fraction of
sampled groups whose rewards are all identical, and GU is its complement.  High
ZVF with low reward indicates cold-start collapse because all completions are
unrewarded and the group-relative advantage is flat.  High ZVF with high reward
indicates saturation because the reward environment is already mostly solved.
Low ZVF indicates that the current policy still samples both successes and
failures, giving the update within-group contrast.  In a 22-run de-duplicated
validation set, a fixed first-five-step rule catches both collapsed runs with no
false positives, but only two positive failures are observed and continuous ZVF
ranking is weak.  The defensible claim is therefore triage: ZVF/GU can support
early decisions to stop, warm-start, densify reward, adjust task difficulty, or
change group size.

Several reward environments in the audit remain lower-tier evidence because
they measure proxies rather than end-to-end capability.  The tool-use reward
scores JSON well-formedness, tool-name selection, and argument-key presence
without executing tools.  The HumanEval-style sandbox executes completions in a
restricted environment with Python built-ins removed, rejecting some legitimate
programs that use functions such as \texttt{len}, \texttt{range}, or
\texttt{sum}.  The MATH reward combines \texttt{\textbackslash boxed\{\}}
format credit with final-answer matching.  In addition, the MATH-500 and
HumanEval prompt pools used inside the training loop are loaded from the
\texttt{test} splits of their respective datasets, so they function here as
reward-environment probes rather than held-out generalization tests.  Positive
rows in these settings are useful evidence about reward dynamics and verifier
interaction, not proof of broad tool use, coding, or mathematical reasoning.

Algorithm comparisons are also lower-tier unless the full stack is fixed or
modeled.  The PPO/GRPO/DPO label alone is not an experimental treatment.  Our
Qwen PPO/GRPO comparison is artifact-sensitive: the source-of-truth ledger
records PPO at 0.225 last-10 reward, while a statistical summary records 0.350.
The Llama-3.1-8B-Instruct comparison favors PPO over the group-relative runner.
These observations should not be converted into a universal PPO-versus-GRPO
ranking.  They instead motivate stack-conditioned reporting of model family,
backend, reward implementation, rollout sampler, optimizer and KL handling,
LoRA target modules, checkpoint-selection rule, and evaluation slice.

\begin{table}[t]
\centering
\caption{Claim--evidence--verdict hierarchy used throughout the paper.}
\label{tab:claim_evidence_verdict}
\small
\begin{tabular}{@{}p{0.27\linewidth}p{0.18\linewidth}p{0.33\linewidth}p{0.16\linewidth}@{}}
\toprule
\textbf{Claim} & \textbf{Status} & \textbf{Headline evidence} & \textbf{Verdict}\\
\midrule
C1: The GRPO-style runner is a reward-contrast amplifier. & Most original contribution &
Mixed-group probability, early-run ZVF/GU, measured base-hit-rate audits, and collapse/saturation cases. &
Signal exists only when sampled groups contain reward diversity.\\
C2: PPO/GRPO/DPO labels are under-specified treatments. & Methodological contribution &
Qwen artifact sensitivity, the Llama PPO-favored row, and backend/model/reward/parser confounding. &
No universal algorithm ranking.\\
C3: Training reward does not equal capability. & Central negative result &
Held-out GSM8K $82.0\%\!\to\!83.3\%$ ($p=0.26$); tool-use rewards are
proxy/schema rewards; HumanEval/MATH harnesses are limited. &
Training reward supports dynamics claims only.\\
\bottomrule
\end{tabular}
\end{table}

The contributions follow this hierarchy.  First, we report the negative
held-out GSM8K result as the boundary on capability claims.  Second, we give a
mechanistic account of when a group-relative runner has a learning signal via
mixed-reward sampling probability.  Third, we introduce ZVF/GU as early-run
triage diagnostics for collapse and saturation.  Fourth, we document how proxy
rewards, parser choices, and sandbox constraints affect the interpretation of
training reward.  Fifth, we provide a source-of-truth ledger and reporting
recommendations showing why PPO, GRPO, and DPO comparisons require full-stack
specification.  We release the code, run logs, and checkpoints at
\url{https://anonymous.4open.science/r/tinker-rl-lab}.

% =============================================================================
\section{Introduction}
\label{sec:intro}

This paper separates three kinds of evidence before making any claim about
reinforcement-learning post-training.  Held-out evaluations are the only basis
we use for capability claims.  Online rewards are evidence about optimization
inside a particular reward environment.  Proxy rewards, partial harnesses, and
stack-specific comparisons are treated as hypotheses or diagnostics unless they
are paired with matched held-out evaluation.  This hierarchy is necessary
because PPO~\citep{schulman2017proximal}, GRPO~\citep{shao2024deepseekmath},
and DPO~\citep{rafailov2024direct} are often discussed as algorithm-level
interventions, while LLM post-training outcomes also depend on base model,
instruction tuning, reward parser, group construction, sampler, backend, LoRA
configuration, checkpoint selection, and evaluation harness.  The resulting
identification problem is the LLM post-training analogue of the reproducibility
concerns identified by Henderson et al.~\citep{henderson2018deep} for deep RL.

On the strongest evidence tier, this paper does \emph{not} show that GRPO
broadly improves reasoning.  The clean held-out capability check is a GSM8K
200-example slice that was not used during training, and it improves only from
$82.0\%$ to $83.3\%$ ($p=0.26$).  We therefore do not claim a significant
held-out math effect.  This negative result determines the scope of the rest of
the paper: training reward curves and proxy environments may explain when a
post-training loop has signal, but they do not establish general reasoning or
tool-use capability.

Within that scope, \ours{} is a systems audit of a group-relative RL runner
rather than an algorithm leaderboard.  The corpus contains 79 runs across
structured tool-call emission, GSM8K, MATH-style tasks, HumanEval-style binary
verification, small and frontier-scale model families, and several training
backends.  The runner is inspired by GRPO, but it is not a faithful reproduction
of canonical GRPO: it keeps group-normalized advantages and a critic-free
structure, but it does not implement the PPO probability ratio and clip of
\citet{schulman2017proximal}, does not keep a frozen reference policy for KL
regularization, performs one optimizer step per rollout batch, and does not use
a completion-only token mask.  A P1-A diagnostic on Qwen3.5-4B measures
$61.6\%$--$89.6\%$ of the full-sequence loss magnitude as coming from prompt
tokens rather than completion tokens.  Claims below therefore concern this
runner and its reward environments, not a universal property of canonical GRPO.

The paper's main positive claim sits one tier below held-out capability: the
runner behaves as a reward-contrast amplifier.  In a critic-free group-relative
update, a prompt is useful only when the sampled group contains distinguishable
reward outcomes.  Under binary or verifier-like rewards, if the current policy
succeeds on prompt $x$ with probability $p_x$ and the group size is $G$, the
probability of obtaining a mixed, gradient-carrying group is approximated by
\[
P(\text{usable}) \approx \frac{1}{N}\sum_x \left[1 - (1-p_x)^G - p_x^G\right].
\]
This mechanism explains why warm-starts, group size, reward density, and task
difficulty matter more than task labels.  The familiar $p_x > 1/G$ rule is only
an expected-one-success heuristic; it is not a threshold theorem for learning.

Zero-Variance Fraction (ZVF) and Gradient Utilization (GU) operationalize that
mechanism as diagnostics, not as capability metrics.  ZVF is the fraction of
sampled groups whose rewards are all identical, and GU is its complement.  High
ZVF with low reward indicates cold-start collapse because all completions are
unrewarded and the group-relative advantage is flat.  High ZVF with high reward
indicates saturation because the reward environment is already mostly solved.
Low ZVF indicates that the current policy still samples both successes and
failures, giving the update within-group contrast.  In a 22-run de-duplicated
validation set, a fixed first-five-step rule catches both collapsed runs with no
false positives, but only two positive failures are observed and continuous ZVF
ranking is weak.  The defensible claim is therefore triage: ZVF/GU can support
early decisions to stop, warm-start, densify reward, adjust task difficulty, or
change group size.

Several reward environments in the audit remain lower-tier evidence because
they measure proxies rather than end-to-end capability.  The tool-use reward
scores JSON well-formedness, tool-name selection, and argument-key presence
without executing tools.  The HumanEval-style sandbox executes completions in a
restricted environment with Python built-ins removed, rejecting some legitimate
programs that use functions such as \texttt{len}, \texttt{range}, or
\texttt{sum}.  The MATH reward combines \texttt{\textbackslash boxed\{\}}
format credit with final-answer matching.  In addition, the MATH-500 and
HumanEval prompt pools used inside the training loop are loaded from the
\texttt{test} splits of their respective datasets, so they function here as
reward-environment probes rather than held-out generalization tests.  Positive
rows in these settings are useful evidence about reward dynamics and verifier
interaction, not proof of broad tool use, coding, or mathematical reasoning.

Algorithm comparisons are also lower-tier unless the full stack is fixed or
modeled.  The PPO/GRPO/DPO label alone is not an experimental treatment.  Our
Qwen PPO/GRPO comparison is artifact-sensitive: the source-of-truth ledger
records PPO at 0.225 last-10 reward, while a statistical summary records 0.350.
The Llama-3.1-8B-Instruct comparison favors PPO over the group-relative runner.
These observations should not be converted into a universal PPO-versus-GRPO
ranking.  They instead motivate stack-conditioned reporting of model family,
backend, reward implementation, rollout sampler, optimizer and KL handling,
LoRA target modules, checkpoint-selection rule, and evaluation slice.

\begin{table}[t]
\centering
\caption{Claim--evidence--verdict hierarchy used throughout the paper.}
\label{tab:claim_evidence_verdict}
\small
\begin{tabular}{@{}p{0.27\linewidth}p{0.18\linewidth}p{0.33\linewidth}p{0.16\linewidth}@{}}
\toprule
\textbf{Claim} & \textbf{Status} & \textbf{Headline evidence} & \textbf{Verdict}\\
\midrule
C1: The GRPO-style runner is a reward-contrast amplifier. & Most original contribution &
Mixed-group probability, early-run ZVF/GU, measured base-hit-rate audits, and collapse/saturation cases. &
Signal exists only when sampled groups contain reward diversity.\\
C2: PPO/GRPO/DPO labels are under-specified treatments. & Methodological contribution &
Qwen artifact sensitivity, the Llama PPO-favored row, and backend/model/reward/parser confounding. &
No universal algorithm ranking.\\
C3: Training reward does not equal capability. & Central negative result &
Held-out GSM8K $82.0\%\!\to\!83.3\%$ ($p=0.26$); tool-use rewards are
proxy/schema rewards; HumanEval/MATH harnesses are limited. &
Training reward supports dynamics claims only.\\
\bottomrule
\end{tabular}
\end{table}

The contributions follow this hierarchy.  First, we report the negative
held-out GSM8K result as the boundary on capability claims.  Second, we give a
mechanistic account of when a group-relative runner has a learning signal via
mixed-reward sampling probability.  Third, we introduce ZVF/GU as early-run
triage diagnostics for collapse and saturation.  Fourth, we document how proxy
rewards, parser choices, and sandbox constraints affect the interpretation of
training reward.  Fifth, we provide a source-of-truth ledger and reporting
recommendations showing why PPO, GRPO, and DPO comparisons require full-stack
specification.  We release the code, run logs, and checkpoints at
\url{https://anonymous.4open.science/r/tinker-rl-lab}.

% =============================================================================
\section{Introduction}
\label{sec:intro}

This paper separates three kinds of evidence before making any claim about
reinforcement-learning post-training.  Held-out evaluations are the only basis
we use for capability claims.  Online rewards are evidence about optimization
inside a particular reward environment.  Proxy rewards, partial harnesses, and
stack-specific comparisons are treated as hypotheses or diagnostics unless they
are paired with matched held-out evaluation.  This hierarchy is necessary
because PPO~\citep{schulman2017proximal}, GRPO~\citep{shao2024deepseekmath},
and DPO~\citep{rafailov2024direct} are often discussed as algorithm-level
interventions, while LLM post-training outcomes also depend on base model,
instruction tuning, reward parser, group construction, sampler, backend, LoRA
configuration, checkpoint selection, and evaluation harness.  The resulting
identification problem is the LLM post-training analogue of the reproducibility
concerns identified by Henderson et al.~\citep{henderson2018deep} for deep RL.

On the strongest evidence tier, this paper does \emph{not} show that GRPO
broadly improves reasoning.  The clean held-out capability check is a GSM8K
200-example slice that was not used during training, and it improves only from
$82.0\%$ to $83.3\%$ ($p=0.26$).  We therefore do not claim a significant
held-out math effect.  This negative result determines the scope of the rest of
the paper: training reward curves and proxy environments may explain when a
post-training loop has signal, but they do not establish general reasoning or
tool-use capability.

Within that scope, \ours{} is a systems audit of a group-relative RL runner
rather than an algorithm leaderboard.  The corpus contains 79 runs across
structured tool-call emission, GSM8K, MATH-style tasks, HumanEval-style binary
verification, small and frontier-scale model families, and several training
backends.  The runner is inspired by GRPO, but it is not a faithful reproduction
of canonical GRPO: it keeps group-normalized advantages and a critic-free
structure, but it does not implement the PPO probability ratio and clip of
\citet{schulman2017proximal}, does not keep a frozen reference policy for KL
regularization, performs one optimizer step per rollout batch, and does not use
a completion-only token mask.  A P1-A diagnostic on Qwen3.5-4B measures
$61.6\%$--$89.6\%$ of the full-sequence loss magnitude as coming from prompt
tokens rather than completion tokens.  Claims below therefore concern this
runner and its reward environments, not a universal property of canonical GRPO.

The paper's main positive claim sits one tier below held-out capability: the
runner behaves as a reward-contrast amplifier.  In a critic-free group-relative
update, a prompt is useful only when the sampled group contains distinguishable
reward outcomes.  Under binary or verifier-like rewards, if the current policy
succeeds on prompt $x$ with probability $p_x$ and the group size is $G$, the
probability of obtaining a mixed, gradient-carrying group is approximated by
\[
P(\text{usable}) \approx \frac{1}{N}\sum_x \left[1 - (1-p_x)^G - p_x^G\right].
\]
This mechanism explains why warm-starts, group size, reward density, and task
difficulty matter more than task labels.  The familiar $p_x > 1/G$ rule is only
an expected-one-success heuristic; it is not a threshold theorem for learning.

Zero-Variance Fraction (ZVF) and Gradient Utilization (GU) operationalize that
mechanism as diagnostics, not as capability metrics.  ZVF is the fraction of
sampled groups whose rewards are all identical, and GU is its complement.  High
ZVF with low reward indicates cold-start collapse because all completions are
unrewarded and the group-relative advantage is flat.  High ZVF with high reward
indicates saturation because the reward environment is already mostly solved.
Low ZVF indicates that the current policy still samples both successes and
failures, giving the update within-group contrast.  In a 22-run de-duplicated
validation set, a fixed first-five-step rule catches both collapsed runs with no
false positives, but only two positive failures are observed and continuous ZVF
ranking is weak.  The defensible claim is therefore triage: ZVF/GU can support
early decisions to stop, warm-start, densify reward, adjust task difficulty, or
change group size.

Several reward environments in the audit remain lower-tier evidence because
they measure proxies rather than end-to-end capability.  The tool-use reward
scores JSON well-formedness, tool-name selection, and argument-key presence
without executing tools.  The HumanEval-style sandbox executes completions in a
restricted environment with Python built-ins removed, rejecting some legitimate
programs that use functions such as \texttt{len}, \texttt{range}, or
\texttt{sum}.  The MATH reward combines \texttt{\textbackslash boxed\{\}}
format credit with final-answer matching.  In addition, the MATH-500 and
HumanEval prompt pools used inside the training loop are loaded from the
\texttt{test} splits of their respective datasets, so they function here as
reward-environment probes rather than held-out generalization tests.  Positive
rows in these settings are useful evidence about reward dynamics and verifier
interaction, not proof of broad tool use, coding, or mathematical reasoning.

Algorithm comparisons are also lower-tier unless the full stack is fixed or
modeled.  The PPO/GRPO/DPO label alone is not an experimental treatment.  Our
Qwen PPO/GRPO comparison is artifact-sensitive: the source-of-truth ledger
records PPO at 0.225 last-10 reward, while a statistical summary records 0.350.
The Llama-3.1-8B-Instruct comparison favors PPO over the group-relative runner.
These observations should not be converted into a universal PPO-versus-GRPO
ranking.  They instead motivate stack-conditioned reporting of model family,
backend, reward implementation, rollout sampler, optimizer and KL handling,
LoRA target modules, checkpoint-selection rule, and evaluation slice.

\begin{table}[t]
\centering
\caption{Claim--evidence--verdict hierarchy used throughout the paper.}
\label{tab:claim_evidence_verdict}
\small
\begin{tabular}{@{}p{0.27\linewidth}p{0.18\linewidth}p{0.33\linewidth}p{0.16\linewidth}@{}}
\toprule
\textbf{Claim} & \textbf{Status} & \textbf{Headline evidence} & \textbf{Verdict}\\
\midrule
C1: The GRPO-style runner is a reward-contrast amplifier. & Most original contribution &
Mixed-group probability, early-run ZVF/GU, measured base-hit-rate audits, and collapse/saturation cases. &
Signal exists only when sampled groups contain reward diversity.\\
C2: PPO/GRPO/DPO labels are under-specified treatments. & Methodological contribution &
Qwen artifact sensitivity, the Llama PPO-favored row, and backend/model/reward/parser confounding. &
No universal algorithm ranking.\\
C3: Training reward does not equal capability. & Central negative result &
Held-out GSM8K $82.0\%\!\to\!83.3\%$ ($p=0.26$); tool-use rewards are
proxy/schema rewards; HumanEval/MATH harnesses are limited. &
Training reward supports dynamics claims only.\\
\bottomrule
\end{tabular}
\end{table}

The contributions follow this hierarchy.  First, we report the negative
held-out GSM8K result as the boundary on capability claims.  Second, we give a
mechanistic account of when a group-relative runner has a learning signal via
mixed-reward sampling probability.  Third, we introduce ZVF/GU as early-run
triage diagnostics for collapse and saturation.  Fourth, we document how proxy
rewards, parser choices, and sandbox constraints affect the interpretation of
training reward.  Fifth, we provide a source-of-truth ledger and reporting
recommendations showing why PPO, GRPO, and DPO comparisons require full-stack
specification.  We release the code, run logs, and checkpoints at
\url{https://anonymous.4open.science/r/tinker-rl-lab}.

% =============================================================================
\section{Introduction}
\label{sec:intro}

This paper separates three kinds of evidence before making any claim about
reinforcement-learning post-training.  Held-out evaluations are the only basis
we use for capability claims.  Online rewards are evidence about optimization
inside a particular reward environment.  Proxy rewards, partial harnesses, and
stack-specific comparisons are treated as hypotheses or diagnostics unless they
are paired with matched held-out evaluation.  This hierarchy is necessary
because PPO~\citep{schulman2017proximal}, GRPO~\citep{shao2024deepseekmath},
and DPO~\citep{rafailov2024direct} are often discussed as algorithm-level
interventions, while LLM post-training outcomes also depend on base model,
instruction tuning, reward parser, group construction, sampler, backend, LoRA
configuration, checkpoint selection, and evaluation harness.  The resulting
identification problem is the LLM post-training analogue of the reproducibility
concerns identified by Henderson et al.~\citep{henderson2018deep} for deep RL.

On the strongest evidence tier, this paper does \emph{not} show that GRPO
broadly improves reasoning.  The clean held-out capability check is a GSM8K
200-example slice that was not used during training, and it improves only from
$82.0\%$ to $83.3\%$ ($p=0.26$).  We therefore do not claim a significant
held-out math effect.  This negative result determines the scope of the rest of
the paper: training reward curves and proxy environments may explain when a
post-training loop has signal, but they do not establish general reasoning or
tool-use capability.

Within that scope, \ours{} is a systems audit of a group-relative RL runner
rather than an algorithm leaderboard.  The corpus contains 79 runs across
structured tool-call emission, GSM8K, MATH-style tasks, HumanEval-style binary
verification, small and frontier-scale model families, and several training
backends.  The runner is inspired by GRPO, but it is not a faithful reproduction
of canonical GRPO: it keeps group-normalized advantages and a critic-free
structure, but it does not implement the PPO probability ratio and clip of
\citet{schulman2017proximal}, does not keep a frozen reference policy for KL
regularization, performs one optimizer step per rollout batch, and does not use
a completion-only token mask.  A P1-A diagnostic on Qwen3.5-4B measures
$61.6\%$--$89.6\%$ of the full-sequence loss magnitude as coming from prompt
tokens rather than completion tokens.  Claims below therefore concern this
runner and its reward environments, not a universal property of canonical GRPO.

The paper's main positive claim sits one tier below held-out capability: the
runner behaves as a reward-contrast amplifier.  In a critic-free group-relative
update, a prompt is useful only when the sampled group contains distinguishable
reward outcomes.  Under binary or verifier-like rewards, if the current policy
succeeds on prompt $x$ with probability $p_x$ and the group size is $G$, the
probability of obtaining a mixed, gradient-carrying group is approximated by
\[
P(\text{usable}) \approx \frac{1}{N}\sum_x \left[1 - (1-p_x)^G - p_x^G\right].
\]
This mechanism explains why warm-starts, group size, reward density, and task
difficulty matter more than task labels.  The familiar $p_x > 1/G$ rule is only
an expected-one-success heuristic; it is not a threshold theorem for learning.

Zero-Variance Fraction (ZVF) and Gradient Utilization (GU) operationalize that
mechanism as diagnostics, not as capability metrics.  ZVF is the fraction of
sampled groups whose rewards are all identical, and GU is its complement.  High
ZVF with low reward indicates cold-start collapse because all completions are
unrewarded and the group-relative advantage is flat.  High ZVF with high reward
indicates saturation because the reward environment is already mostly solved.
Low ZVF indicates that the current policy still samples both successes and
failures, giving the update within-group contrast.  In a 22-run de-duplicated
validation set, a fixed first-five-step rule catches both collapsed runs with no
false positives, but only two positive failures are observed and continuous ZVF
ranking is weak.  The defensible claim is therefore triage: ZVF/GU can support
early decisions to stop, warm-start, densify reward, adjust task difficulty, or
change group size.

Several reward environments in the audit remain lower-tier evidence because
they measure proxies rather than end-to-end capability.  The tool-use reward
scores JSON well-formedness, tool-name selection, and argument-key presence
without executing tools.  The HumanEval-style sandbox executes completions in a
restricted environment with Python built-ins removed, rejecting some legitimate
programs that use functions such as \texttt{len}, \texttt{range}, or
\texttt{sum}.  The MATH reward combines \texttt{\textbackslash boxed\{\}}
format credit with final-answer matching.  In addition, the MATH-500 and
HumanEval prompt pools used inside the training loop are loaded from the
\texttt{test} splits of their respective datasets, so they function here as
reward-environment probes rather than held-out generalization tests.  Positive
rows in these settings are useful evidence about reward dynamics and verifier
interaction, not proof of broad tool use, coding, or mathematical reasoning.

Algorithm comparisons are also lower-tier unless the full stack is fixed or
modeled.  The PPO/GRPO/DPO label alone is not an experimental treatment.  Our
Qwen PPO/GRPO comparison is artifact-sensitive: the source-of-truth ledger
records PPO at 0.225 last-10 reward, while a statistical summary records 0.350.
The Llama-3.1-8B-Instruct comparison favors PPO over the group-relative runner.
These observations should not be converted into a universal PPO-versus-GRPO
ranking.  They instead motivate stack-conditioned reporting of model family,
backend, reward implementation, rollout sampler, optimizer and KL handling,
LoRA target modules, checkpoint-selection rule, and evaluation slice.

\begin{table}[t]
\centering
\caption{Claim--evidence--verdict hierarchy used throughout the paper.}
\label{tab:claim_evidence_verdict}
\small
\begin{tabular}{@{}p{0.27\linewidth}p{0.18\linewidth}p{0.33\linewidth}p{0.16\linewidth}@{}}
\toprule
\textbf{Claim} & \textbf{Status} & \textbf{Headline evidence} & \textbf{Verdict}\\
\midrule
C1: The GRPO-style runner is a reward-contrast amplifier. & Most original contribution &
Mixed-group probability, early-run ZVF/GU, measured base-hit-rate audits, and collapse/saturation cases. &
Signal exists only when sampled groups contain reward diversity.\\
C2: PPO/GRPO/DPO labels are under-specified treatments. & Methodological contribution &
Qwen artifact sensitivity, the Llama PPO-favored row, and backend/model/reward/parser confounding. &
No universal algorithm ranking.\\
C3: Training reward does not equal capability. & Central negative result &
Held-out GSM8K $82.0\%\!\to\!83.3\%$ ($p=0.26$); tool-use rewards are
proxy/schema rewards; HumanEval/MATH harnesses are limited. &
Training reward supports dynamics claims only.\\
\bottomrule
\end{tabular}
\end{table}

The contributions follow this hierarchy.  First, we report the negative
held-out GSM8K result as the boundary on capability claims.  Second, we give a
mechanistic account of when a group-relative runner has a learning signal via
mixed-reward sampling probability.  Third, we introduce ZVF/GU as early-run
triage diagnostics for collapse and saturation.  Fourth, we document how proxy
rewards, parser choices, and sandbox constraints affect the interpretation of
training reward.  Fifth, we provide a source-of-truth ledger and reporting
recommendations showing why PPO, GRPO, and DPO comparisons require full-stack
specification.  We release the code, run logs, and checkpoints at
\url{https://anonymous.4open.science/r/tinker-rl-lab}.
